%============================================================ 1. 서론
\section{서론}
\label{sec:intro}

\subsection{배경}

다수의 저비용 무인수상정(Unmanned Surface Vehicle, USV)이 고가치 함정 한 척을 향해 여러
방향에서 동시에 돌진하는 자폭 공격은 대표적인 비대칭 위협이다. 유한한 방어 자산을 다수
표적에 배분하는 문제는 무기-표적 할당(weapon-target assignment)으로 오래 연구되어
왔으나~\citep{ahuja2007wta}, 표적 수가 요격 자산을 넘어서면 최적 할당으로도 처리되지 않는
표적이 남는다.

무인 위협에 대한 기존 대응은 두 갈래로 구분된다~\citep{wang2021cuas}. 표적의 통신·항법
·센서를 교란하는 연성 대응(soft-kill)은 위협 자체를 제거하지 못하며, 전파를 방사하지 않는
유선 조종 표적에는 작용하지 않는다~\citep{sutea2026fiber}. 표적을 직접 타격하는 경성
대응(hard-kill)은 위협을 확실히 제거하지만, 수천 달러 규모의 표적에 발당 수십만에서
수백만 달러의 요격 유도탄을 소모하는 비용 구조를 지탱하기 어렵고~\citep{rumbaugh2024cost},
한 번에 한 표적만 처리하므로 표적 수가 교전 용량을 넘어서면 포화가 발생한다.

본 연구는 대안으로 저비용 그물을 이용한 물리적 포획을 채택한다. 이는 무인기 대응 분야에서
독립된 대응 범주로 다루어져 왔으며~\citep{wang2021cuas}, 표적의 통신 방식과 무관하게
작용하고 표적을 정지시킨다는 점에서 앞의 두 한계를 비켜간다. 방어정은 적선보다 느리므로
추격으로 따라잡을 수 없고, 대신 적선이 지나갈 접근 회랑을 앞질러 그물벽으로 차단한다
(Fig.~\ref{fig:concept}).

%-------------------------------------------------- 그림 1
% 2단에서 단폭(71.7 mm)에 넣으면 HUD 글자가 2.5 pt 로 판독 불가다. 양단 걸침으로
% 170 mm 를 확보해 약 5.9 pt 로 올린다 --- 캡션이 독자에게 이 화면을 읽으라고 요구한다.
\begin{figure*}[t]
\centering
\includegraphics[width=0.86\textwidth]{fig_concept_art.png}
\caption{교전 개념. 중앙의 모선을 향해 자폭 USV가 여러 접근 회랑(붉은 부채꼴)으로
         동시에 돌진한다. 방어정은 적선보다 \emph{느리므로}(Table~\ref{tab:scenario})
         추격이 성립하지 않고, 대신 요격 환형 위에서 회랑을 가로지르는 그물벽을 미리 친다.
         그물벽은 두 지점을 잇는 선분이며, 적선이 여기 얽히면 포획된다. 안쪽 원은 돌파
         판정 반경으로, 적선이 이 안에 들어오면 방어 실패로 계수한다.}
\label{fig:concept}
\end{figure*}

\subsection{문제의식}

이 방어 문제는 성격이 다른 두 판단을 동시에 요구한다. 하나는 \emph{어느 방어정이 어느
적 무리를 맡을 것인가}라는 이산적·전술적 배정이고, 다른 하나는 \emph{맡은 무리를 막기 위해
해면의 어느 지점에 그물벽을 놓을 것인가}라는 연속적·기하적 기동이다. 전자는 전장 전체의
맥락과 우선순위 판단을 요구하고, 후자는 국소적이지만 매 결정마다 새로 계산되어야 한다.
두 판단은 같은 주기로 갱신되지만 \emph{산출에 걸리는 시간이 크게 다르다} --- 이 비대칭이
두 계층을 나누는 출발점이다.

본 연구는 느린 쪽인 배정을 언어모델 지휘관에게, 빠른 쪽인 기동을 학습된 정책에 맡긴다.
이때 지휘관의 출력을 조합 최적화로 보정하지 않는다. 보정을 두면 최종 배정이 사실상 비용
함수의 산물이 되어, 성능이 지휘관의 판단에서 온 것인지 후처리에서 온 것인지 가릴 수 없기
때문이다. 이 선택은 언어모델 계층에 불리할 수 있으나, 두 계층의 기여를 \emph{따로 잴 수
있게} 한다.

\subsection{배치를 전제로 한 설계}

본 연구는 시뮬레이션 안에서만 성립하는 체계를 목표로 하지 않았다. 방법을 고를 때마다
\emph{온보드에서 돌 수 있는가}를 함께 물었고, 그 선택들이 본문 곳곳에 흩어져 있으므로
여기에 모아 둔다.

\begin{itemize}[leftmargin=*, itemsep=2pt, topsep=3pt]
\item \textbf{실해역 지형에서 학습·평가한다.} 교전 해역은 가상의 빈 바다가 아니라
      OpenStreetMap 의 실제 해안선(통영 매물도 인근)을 래스터화한 것이며, 육지는 충돌
      판정과 관측 채널 양쪽에 들어간다(\S\ref{sec:obs}).
\item \textbf{지연을 가정하지 않고 측정한다.} 언어모델 지휘관의 응답 지연을 단일
      GPU(RTX 4070) 구성에서 호출 13{,}547회에 걸쳐 실측하였고, 결정 주기를 초과하는
      꼬리가 얼마나 되는지까지 보고한다(\S\ref{sec:res_latency}). 비동기 분리는 그 측정에서 나온 설계다.
\item \textbf{학습과 추론의 동작이 갈리지 않는다.} 배치 통계에 의존하는 정규화를 일절
      쓰지 않으므로(그룹 정규화만 사용) 배치 크기 1의 실시간 추론이 학습과 같은 연산을
      수행한다. 통계 워밍업이 필요 없다(\S\ref{sec:deploy}).
\item \textbf{출력이 곧 항법 명령이다.} 정책의 출력은 제어 신호가 아니라 경유점 두 개이며,
      이는 ROS2 항법 스택이 그대로 받는 형식이다. 저수준 추종은 학습과 분리되어 있어
      실제 선체의 동역학으로 교체할 수 있다(\S\ref{sec:deploy}).
\end{itemize}

\subsection{기여}

본 논문의 기여는 \emph{무엇이 새로 가능해졌는가}로 정리된다.

\begin{enumerate}[leftmargin=*, itemsep=3pt, topsep=4pt]
\item \textbf{적보다 느린 소수의 방어정으로도 군집을 물리적으로 막을 수 있다 --- 그것도
      훨씬 적은 자원으로.} 그물벽을 놓을 자리를 해면 전체에 대한 점수맵 위에서 고르게 하여,
      같은 정도로 막으면서 포획당 그물 소모를 $44.5\,\%$, 총 선회량을 $46.6\,\%$ 줄였다.
      \textbf{이 가운데 언어모델 없이 기동 계층 단독으로 얻는 몫}은 포획당 그물 $22.9\,\%$,
      그물 걸림 $24.3\,\%$ 감소이다.
      기법 쪽 근거는 15채널 격자 관측과 FiLM 조건화 U-Net, 그리고 연속 좌표
      회귀 대신 픽셀을 지목하는 행동 표현이다.
\item \textbf{방어정 수와 적선 수가 바뀌어도 같은 모델이 돈다.} 관측을 선박별
      벡터의 나열이 아니라 \emph{이미지 형식}의 15채널 격자로 바꾼 것이 핵심이다.
      전장의 선박 수가 몇 척이든 관측은 같은 크기의 격자 위에 그려지므로 입력의 형태가
      유지되고, 입력 순열 문제 자체가 생기지 않는다. 가변 개수 입력을 다루는 통상의
      방법 --- 최대 개수로 패딩하고 마스크를 붙이거나, 개체 목록 위에 어텐션을 얹는
      방식 --- 과 달리, 여기서는 개체 수에 따라 달라지는 것이 아무것도 없다. 계산량이
      선박 수에 비례해 늘지 않고, 마스크나 최대 개수 같은 설계 상수를 둘 필요가 없으며,
      개체로 표현되지 않는 정보(해안선·섬 같은 육지, 모선 주변의 위협 분포)도 같은
      격자에 채널 하나로 얹힌다. 무엇보다 출력이 같은 격자 위의 픽셀이므로 관측과 행동이
      하나의 좌표계를 공유한다. 여기에 방어정 간 파라미터 공유를 더해 모델 크기 역시
      방어정 수에 불변이다(총 244{,}389개). 그 결과 편성이 달라질 때마다 모델을 다시
      만들 필요 없이, 동일한 하나의 모델이 어떤 편성·어떤 상황에서도 강건하게 대응한다.
\item \textbf{협동을 보상 설계가 아니라 크레딧 배분으로 얻었다.} 팀 보상을 모든 방어정이
      그대로 나누어 가지면 잘한 배와 못한 배가 같은 신호를 받아 협동이 학습되지 않는다.
      본 연구는 그룹 상대 정책 최적화를 다중 방어정으로 확장하고, 한 척의 행동만 바꾼 채
      시뮬레이터를 실제로 재진행하여 얻은 차이를 그 배의 몫으로 돌리는 반사실적
      크레딧으로 개별 기여를 분리한다. 학습된 비평가로 주변화하지 않으므로 가치 추정
      편향과 비평가 학습의 불안정성이 설계에서 사라진다. 그 결과 협동을 유도하는 보상
      항을 손으로 설계하지 않고도, 하나의 팀 보상만으로 다중 방어정의 협동 기동이
      학습된다.
\item \textbf{언어모델을 지휘 계층에 넣을 때의 판단 근거를 수치로 제공한다.} 배정 권한을
      지휘관에게 온전히 넘기고 조합 최적화 후처리를 제거한 구성에서, 8조건 $\times$
      3포메이션 $\times$ 30시드(720 에피소드)를 시드 대응표본으로 수집해 두 계층의 기여를
      분리 측정하였다. \textbf{언어모델이 항상 이득은 아니다} --- 코드 휴리스틱을 넘어서는
      능력 임계가 존재하며, 그 아래 모델은 오히려 성능을 떨어뜨린다.
\item \textbf{지휘 계층을 무엇으로 평가해야 하는지 밝혔다.} 성능을 설명하는 것은 배정
      커버리지가 아니라 \emph{재계획 간 계획 안정성}이다. 가장 많이 배정한 모델이 가장
      나빴다. 주기적으로 재계획하는 구조에서는 매 주기의 국소 최적성보다 시간에 걸친
      일관성이 중요하다는 설계 원리가 여기서 나온다.
\end{enumerate}

\noindent
논문의 구성은 다음과 같다. \S\ref{sec:related}에서 관련 연구를 정리하고,
\S\ref{sec:problem}에서 문제와 평가 지표를 정의한다. \S\ref{sec:method}에서 제안 체계를
기술하고, \S\ref{sec:experiment}에서 실험 설계와 통계 절차를 밝힌다.
\S\ref{sec:results}에서 결과를, \S\ref{sec:discussion}에서 해석을 제시하고,
\S\ref{sec:limitations}에서 한계를 논의한다.
